\documentclass[aspectratio=169,12pt]{beamer}
\usepackage{fontspec}
\usepackage[portuguese]{babel}
\usepackage{xcolor,listings,tikz,booktabs,amsmath,amssymb,graphicx,multicol,array,colortbl}
\usetikzlibrary{arrows.meta,positioning,shapes.geometric,calc}

\definecolor{navy}{HTML}{0D1B3E}
\definecolor{gold}{HTML}{F5A623}
\definecolor{qgreen}{HTML}{1E8B4C}
\definecolor{qblue}{HTML}{1A6EAE}
\definecolor{codebg}{HTML}{EEF3F8}

\usetheme{default}
\useinnertheme{circles}
\setbeamertemplate{navigation symbols}{}
\setbeamertemplate{footline}{%
  \hfill{\color{gray!60}\scriptsize\insertframenumber/\inserttotalframenumber}\hspace{.6em}\vspace{.4em}}

\setbeamercolor{frametitle}{bg=navy,fg=white}
\setbeamercolor{structure}{fg=navy}
\setbeamercolor{alerted text}{fg=gold}
\setbeamercolor{item}{fg=gold}
\setbeamercolor{subitem}{fg=navy}
\setbeamercolor{block title}{bg=qblue,fg=white}
\setbeamercolor{block body}{bg=qblue!8,fg=black}
\setbeamercolor{block title alerted}{bg=gold,fg=navy}
\setbeamercolor{block body alerted}{bg=gold!15,fg=navy}
\setbeamercolor{block title example}{bg=qgreen,fg=white}
\setbeamercolor{block body example}{bg=qgreen!10,fg=black}
\setbeamerfont{frametitle}{size=\large,series=\bfseries}
\setbeamerfont{block title}{series=\bfseries}

\setbeamertemplate{title page}{%
  \begin{tikzpicture}[remember picture,overlay]
    \fill[navy](current page.north west)rectangle(current page.south east);
    \fill[gold]([yshift=.7cm]current page.south west)rectangle(current page.south east);
    \node[anchor=south,yshift=.73cm,navy,font=\tiny\bfseries]at(current page.south)
      {INTENSIVÃO QUANT AI \textbullet{} IMPACT UFSCAR};
  \end{tikzpicture}
  \vspace{.65cm}
  \begin{center}
    {\color{gold}\scriptsize\bfseries INTENSIVÃO QUANT AI \textbullet{} IMPACT UFSCAR}\\[.4cm]
    {\color{white}\LARGE\bfseries\inserttitle}\\[.25cm]
    {\color{gold!85}\normalsize\insertsubtitle}\\[.5cm]
    {\color{white!70}\small\insertauthor}\\[.1cm]
    {\color{white!50}\footnotesize\insertdate}
  \end{center}}

\lstdefinestyle{python}{language=Python,
  basicstyle=\ttfamily\scriptsize\color{qblue},
  keywordstyle=\color{navy}\bfseries,
  stringstyle=\color{qgreen!80!black},
  commentstyle=\color{gray!70}\itshape,
  backgroundcolor=\color{codebg},
  frame=l,framerule=2pt,rulecolor=\color{gold},
  xleftmargin=1em,breaklines=true,showstringspaces=false,tabsize=4,
  extendedchars=false,inputencoding=ascii}
\lstset{style=python}

\author{Felipe Maldonado\\{\scriptsize VP Diretoria Quant~\textbullet{}~Impact UFSCAR}}
\date{Intensivão Quant AI 2025}
\title{Aula 07 --- Sinal v2}
\subtitle{Vol-Adjusted Momentum: Melhorando o Sinal}

\begin{document}

\begin{frame}[plain]
\titlepage
\end{frame}

\begin{frame}[fragile]{Nesta Aula}
\begin{columns}[T]
  \column{0.5\textwidth}
\begin{block}{Teoria (30 min)}
\begin{itemize}
  \item Limitação crítica do sinal v1: momentum crash
  \item Vol-adjusted momentum: Moskowitz, Ooi \& Pedersen (2012)
  \item Comparação fair entre sinais: IC médio e Sharpe OOS
  \item Quando o ajuste de volatilidade ajuda (e quando não)
\end{itemize}
\end{block}
  \column{0.47\textwidth}
\begin{block}{Código (70 min)}
\begin{itemize}
  \item \texttt{calcular\_vol\_rolling(ret, janela=63)} — vol diária rolling
  \item \texttt{calcular\_sinal\_v2(sinal\_v1, vol)} — sinal normalizado
  \item Comparação de IC: v1 vs v2 (bar chart)
  \item Backtest comparativo: equity curves lado a lado
  \item Salvar \texttt{sinal\_v2.parquet} e \texttt{pesos\_v2\_final.parquet}
\end{itemize}
\end{block}
\end{columns}
\end{frame}

\begin{frame}[fragile]{Limitação do Sinal v1 — Momentum Crash}
\begin{columns}[T]
  \column{0.5\textwidth}
\textbf{Problema identificado por Daniel \& Moskowitz (2016):}
\vspace{.3em}
\begin{itemize}
  \item Em crises de mercado (2008, 2020), os ``vencedores'' do sinal v1 são tipicamente ativos de \textbf{alta beta}
  \item No crash, estes ativos caem mais que o mercado
  \item O sinal de momentum ``captura'' um risco de mercado latente
  \item Resultado: drawdown de 40--60\% em crises severas
\end{itemize}
\vspace{.3em}
\begin{alertblock}{Sinal 12-1 sem ajuste de vol}
  Compra ativos de alto retorno recente $\approx$ compra ativos de alta vol/beta.
  Em crises, isso amplifica as perdas.
\end{alertblock}
  \column{0.47\textwidth}
\begin{block}{Solução: Vol-Adjustment}

\textbf{Moskowitz, Ooi \& Pedersen (2012):}
\[
\text{Sinal}^{v2}_{i,t} = \frac{\text{Sinal}^{v1}_{i,t}}{\hat\sigma_{i,t}}
\]
$\hat\sigma_{i,t}$: volatilidade rolling dos últimos 63 dias úteis.
\vspace{.3em}
\begin{itemize}
  \item Normaliza o sinal pela volatilidade corrente
  \item Ativos mais voláteis recebem sinal ``deflacionado''
  \item Melhora IC e reduz tail risk
  \item Alinha com risk parity dos pesos
\end{itemize}
\end{block}\vspace{.3em}
\begin{exampleblock}{Evidência}
Moskowitz et al.\ documentam melhora de Sharpe em ações, bonds, commodities e moedas.
\end{exampleblock}
\end{columns}
\end{frame}

\begin{frame}[fragile]{Comparação Fair: v1 vs v2}
\begin{columns}[T]
  \column{0.5\textwidth}
\textbf{Protocolo de comparação:}
\begin{enumerate}
  \item Mesmo universo de ativos
  \item Mesmo período de análise
  \item Mesma regra de portfólio (top 20\%, equal-weight)
  \item Mesma modelagem de custos (turnover $\times$ 0{,}1\%)
\end{enumerate}
\vspace{.3em}
\textbf{Métricas de comparação:}
\begin{itemize}
  \item IC médio: poder preditivo bruto do sinal
  \item Sharpe ratio do portfólio
  \item Maximum drawdown
  \item Correlation entre sinal v1 e v2: quanto diferem?
\end{itemize}
  \column{0.47\textwidth}
\begin{block}{Quando v2 ajuda?}
\begin{itemize}
  \item Ambientes de alta volatilidade (crises, pandemia)
  \item Quando os ``vencedores'' são concentrados em setores de alto beta
  \item Quando há momentum crash latente (beta do portfólio > 1{,}2)
\end{itemize}
\end{block}\vspace{.3em}
\begin{alertblock}{Quando v2 não ajuda?}
\begin{itemize}
  \item Mercados trending com baixa dispersão de vol
  \item Quando todos os ativos têm vol similar (correlação alta)
  \item Períodos de vol muito baixa — normalização é instável
\end{itemize}
\end{alertblock}
\end{columns}
\end{frame}

\begin{frame}[fragile]{O que Vamos Construir}
\begin{columns}[T]
  \column{0.52\textwidth}
\begin{block}{Funções a implementar}
\begin{itemize}
  \item \texttt{calcular\_vol\_rolling(ret\_diarios, janela=63)} $\to$ vol mensal estimada
  \item \texttt{alinhar\_vol\_com\_sinal(vol, sinal)} $\to$ vol no mesmo index
  \item \texttt{calcular\_sinal\_v2(sinal\_v1, vol)} $\to$ sinal normalizado
  \item \texttt{comparar\_ic(sinal\_v1, sinal\_v2, ret)} $\to$ DataFrame comparativo
\end{itemize}
\end{block}
  \column{0.45\textwidth}
\begin{block}{Entradas (parquets)}
\begin{itemize}
  \item \texttt{sinal\_v1.parquet}
  \item \texttt{retornos\_diarios\_limpo.parquet}
  \item \texttt{retornos\_mensais\_limpo.parquet}
\end{itemize}
\end{block}
\vspace{.4em}
\begin{exampleblock}{Saídas (parquets)}
\begin{itemize}
  \item \texttt{sinal\_v2.parquet}
  \item \texttt{pesos\_v2.parquet}
  \item Gráfico comparativo IS v1 vs v2
\end{itemize}
\end{exampleblock}
\end{columns}
\end{frame}

\begin{frame}[fragile]{Referências}
\begin{multicols}{2}
\tiny
\begin{itemize}
  \item \textbf{Moskowitz, T., Ooi, Y. \& Pedersen, L.} (2012). Time series momentum. \textit{Journal of Financial Economics}, 104(2), 228--250.
  \item \textbf{Daniel, K. \& Moskowitz, T.} (2016). Momentum crashes. \textit{JFE}, 122(2), 221--247.
  \item \textbf{Barroso, P. \& Santa-Clara, P.} (2015). Momentum has its moments. \textit{JFE}, 116(1), 111--120.
  \item \textbf{Jegadeesh, N. \& Titman, S.} (1993). Returns to buying winners. \textit{Journal of Finance}, 48(1), 65--91.
  \item \textbf{Asness, C. et al.} (2013). Value and momentum everywhere. \textit{Journal of Finance}, 68(3), 929--985.
\end{itemize}
\end{multicols}
\end{frame}


\end{document}
